% =====================================================================
%  R. Nielsen, PROPÆDEUTIK OG LOGIK.  Cursus for Universitetsaaret 1868-69.
%  Kjøbenhavn: Forlagt af den Gyldendalske Boghandel (F. Hegel); trykt hos J. H. Schultz, 1868.  6 + 456 pp.  Fraktur.
%
%  THE NEW PAGES ONLY.  The volume is a composite: 445 of its 456 text pages are sheets of three earlier books in the
%  same setting (the same type, line for line), and those books are transcribed in this repository.  This file holds
%  what was set for 1868: the title page, the separately paginated Oversigt (pp. [1]-6), pp. 81-89 (the opening of
%  section B, »Propædeutik og Ideernes Logik«; pp. 81 and 89 given whole although their first 17 / last 25 lines are the
%  earlier setting) and pp. 263-264 (the seam into section C, »Propædeutik og Ideelære«).  The rest:
%    pp. 1-80 (and p. 81, lines 1-17) = Forelæsninger over »Philosophisk Propædeutik« 1860-61 (1862), pp. 1-81
%                                       (texts/nielsen/propaedeutik-1860-61)
%    pp. 90-262 (and p. 89 from »1) Ordet og Gud«) = Grundideernes Logik II (1866), the same pages
%                                       (texts/nielsen/grundideernes-logik-2)
%    pp. 265-456 = Grundideernes Logik I (1864), the same pages (texts/nielsen/grundideernes-logik)
%  How this was established (map1868.py, concord.py, setting.py, changes.py, concordance.tsv in
%  ebook-method/propaedeutik-1868/): the 1868 scan's text layer matched page for page against the transcriptions, and
%  its sequence of line-end words against the earlier scans' layers (80-100 % on shared pages, 0-9 % on new ones); the
%  earlier books' printed errata are not corrected in 1868, and a three-way word diff found no textual change.
%
%  TRANSCRIPTION -- GENERATED; DO NOT EDIT BY HAND.  Rebuild, chained (&&), from the repo root:
%    python3 ebook-method/propaedeutik-1868/draft.py --emit --force
%    python3 ebook-method/collate/check.py propaedeutik-1868 --witness ocr --apply
%    python3 ebook-method/propaedeutik-1868/oepass.py --apply
%    python3 ebook-method/propaedeutik-1868/finish.py --apply
%
%  WORD-ACCURACY (2026-10-13).  No e-book and no other transcription of these pages exists, except that p. 89 from
%  »1) Ordet og Gud« and p. 263 are the type of Grundideernes Logik II, whose transcription served as their draft (the
%  layer garbles the Greek and interleaves the note).  The draft (Google's layer) was collated word by word, with
%  punctuation and word division, against tesseract (tessdata_best script/Fraktur), with deu_latf and frak2021 as third
%  and fourth readings: 613 disagreements.  100 settled by check.py's own rules, 358 by rules.py (dash+ and case rules
%  withdrawn: Google's 600-ppi scans make spurious dashes), 26 Greek words on p. 89 by hand (GL II's type), 142 at the
%  image in BLIND crops read by two Sonnet readers (controls 4 of 4 right; a blind audit of 32 rule decisions found no
%  rule error, 3 reader slips), 34 by hand from the readers' notes and 4 line crops.  OPEN 0.
%
%  FULL PAGE CHECK (2026-10-13, pagecheck.py, every page: the title leaf and pp. 1-6, 81-89, 263-264, 4613 words).
%  Every page was transcribed BLIND at the image by three further Sonnet readers and compared word by word; every
%  difference was settled in the main session, the doubtful ones in line crops (.render/propaedeutik-1868-adj*).
%  Transcription errors found and mended: 37 of words (8.0 per 1000), almost all of the classes that the layer and
%  tesseract SHARE in Fraktur (k read f: iffe, funne, frafjende; s read f: faadant, fingulære; R read N: Net; V read B:
%  Basen, Bærende; Ø read D: Die, Diet, Diemed; I read J; æ read a: vare, Vasen, formales; Kant read Kaut); 6 of
%  punctuation (two dashes, a parenthesis, three quotation marks); 7 word divisions that the crop readers' line-break
%  forms had brought in (Qva-litet, Gjen-/standen ...); and the emphasis: 42 letterspaced runs, a bold phrase and the
%  bold heads, none of which the OCRs see.  Since every page was read, these are all the errors the check could find;
%  what it cannot see is a reading in which the transcription and the blind reader agree on the same fault.  The
%  readers' Qu for the print's Qv, Character for Characteer, Enhed for Eenhed and the like were their slips (crops).
%  Greek accents on p. 89 follow GL II (the same type); a reader read acutes where GL II has graves, below what the
%  crops resolve.
%  PRINTED AS IS: „(Uende-/hedsbegrebet)“ p. 263 (a syllable lost when the note's first line was reset for 1868).
%  The ø rule (CLAUDE.md): 13 o-forms set as ø (følgende, Størrelse, føies, Dømmende ...) and „løse“ (p. 88).
%
%  PAGE MAP.  Scan ~/bibliotek/Nielsen, Rasmus/Propædeutik-1868.pdf (Google Books, 482 PDF pp.; SCANS.tsv).
%  PDF 9 and 11: the title leaf, scanned twice; PDF 13-18: Oversigt pp. [1]-6 (printed p = PDF - 12); text pp. 1-456 = PDF 19-474
%  (printed p = PDF - 18).  PDF 476: the publisher's advertisement of Nielsen's writings (not transcribed, Hans
%  2026-10-13).  Other PDF pages blank, or Google's notice (PDF 1).
%
% =====================================================================
\documentclass[12pt,a4paper,openany]{book}
\usepackage[utf8]{inputenc}
\usepackage[T1]{fontenc}
\usepackage{libertinus}
\usepackage{graphicx}
\DeclareUnicodeCharacter{0254}{\reflectbox{c}}  % ɔ (U+0254) = reversed c
\DeclareUnicodeCharacter{A75B}{r}   % r rotunda (ꝛc. = etc.)
\DeclareUnicodeCharacter{221E}{\ensuremath{\infty}}  % ∞ (Oversigt p. 1)
\usepackage[danish]{babel}
\usepackage{textalpha}   % polytonic Greek typed directly (p. 89), as in grundideernes-logik-2
\usepackage[protrusion=true,expansion=true]{microtype}
\usepackage{setspace}
\usepackage[margin=1.3in]{geometry}
\usepackage{fancyhdr}
\usepackage{hyperref}
\hypersetup{pdftitle={Propædeutik og Logik. Cursus for Universitetsaaret 1868--69 (the new pages)}, pdfauthor={Rasmus Nielsen}, hidelinks}
\setlength{\headheight}{15pt}
\pagestyle{fancy}
\fancyhf{}
\fancyhead[LE,RO]{\thepage}
\fancyhead[C]{\textit{R. Nielsen: Propædeutik og Logik (1868)}}
\renewcommand{\thefootnote}{\ifcase\value{footnote}\or *)\or **)\or ***)\else
  \arabic{footnote})\fi}
\setstretch{1.3}
\usepackage{marginnote}
\newcommand{\opage}[1]{\marginnote{\footnotesize\textit{[#1]}}}
\newcommand{\apage}[1]{\marginnote{\footnotesize\textit{[#1]}}}
\newcommand{\gr}[1]{#1}   % marks the Greek (p. 89) for tools; textalpha sets it
% the end of a note begun on a page not transcribed here (p. 263: the note of GL II p. 262), set at the foot without a mark
\newcommand{\contnote}[1]{\begingroup\renewcommand{\thefootnote}{}\footnotetext{\textit{[continued from p.~262:]} #1}\endgroup}
% editorial seams: where the 1868 volume continues in the type of an earlier book (no-argument macros, so that the
% collation, which reads only the printed words, does not see them)
\newcommand{\editorial}[1]{\par\medskip\begin{center}\small\textit{[#1]}\end{center}\medskip\par}
\newcommand{\seamOversigt}{\editorial{End of the Oversigt.  The text pp.~1--80, and p.~81 to the head »B.«, is
  \emph{Forelæsninger over »Philosophisk Propædeutik« fra Universitetsaaret 1860--61} (1862), pp.~1--81, in the same
  setting: see that transcription.}}
\newcommand{\seamB}{\editorial{pp.~90--262 (and p.~89 from »1) Ordet og Gud«) are \emph{Grundideernes Logik II} (1866),
  the same pages in the same setting: see that transcription.  p.~263 below is GL II p.~263 with one line taken up
  from GL II p.~264.}}
\newcommand{\seamC}{\editorial{pp.~265--456 are \emph{Grundideernes Logik I} (1864), the same pages in the same
  setting: see that transcription.  End of the volume.}}

\begin{document}



% --- PDF 9, the title page; PDF 11 is the same leaf scanned a second time (cleaner; a pencilled underline under the
% author's name and two strokes).  Read at the image by the page reader: the title large, the course line, the author and
% the place in bold Fraktur; the publisher and printer small; the year in antiqua.
\begin{titlepage}
\begin{center}
\vspace*{4em}
{\Huge\textbf{Propædeutik og Logik.}}\\[2em]
{\Large\textbf{Cursus for Universitetsaaret 1868--69.}}\\[3em]
{\small Af}\\[1em]
{\Large\textbf{R. Nielsen.}}\\[8em]
{\small\textbf{Kjøbenhavn.}}\\[0.5em]
{\small Forlagt af den Gyldendalske Boghandel (F. Hegel.)}\\[0.5em]
{\footnotesize Trykt hos J. H. Schultz.}\\[0.5em]
\textit{1868.}
\end{center}
\end{titlepage}

\chapter*{Editorial note}
\addcontentsline{toc}{chapter}{Editorial note}
{\small This volume, Nielsen's course book for 1868--69, was made up for the most part of sheets of his earlier books,
reissued in the same setting.  Only the pages composed for 1868 are transcribed here; the others are in the
transcriptions of the earlier books:
\begin{center}
\begin{tabular}{ll}
1868 & is \\ \hline
Oversigt pp.~[1]--6 & new (below) \\
pp.~1--80, p.~81 lines 1--17 & \emph{Forelæsninger over »Philosophisk Propædeutik« 1860--61}, pp.~1--81 \\
p.~81 from »B.«, pp.~82--88, p.~89 lines 1--8 & new (below; pp.~81 and 89 given whole) \\
pp.~89 (from »1) Ordet og Gud«)--262 & \emph{Grundideernes Logik II} (1866), the same pages \\
pp.~263--264 & the seam into section C (below; p.~264 is new) \\
pp.~265--456 & \emph{Grundideernes Logik I} (1864), the same pages \\
\end{tabular}
\end{center}
Page numbers in the margin are those of the 1868 print.\par}

% --- p. 1 ---
\setcounter{footnote}{0}%
\opage{1}\begin{center}{\large Oversigt.}\end{center}


Vi aabne os et Indblik i den philosophiske Videns
Eiendommelighed, idet vi begynde med at opfatte Philosophien
som Modsigelseslære, som Mulighedslære og som
Dialektik. For at charakterisere Dialektikens Forhold til de
exacte Analyser tilføies et Overblik over Forholdet mellem den
mathematiske og den philosophiske Behandling af Uendelighedsproblemet.
Af dette Overblik fremhæves følgende Hovedpunkter.

a) Enhver endelig Størrelse kan deles i uendelig mange
uendelig smaa Dele, og Summen angives ved 0.∞. Da
nu den ene endelige Størrelse ikke desto mindre er forskjellig
fra den anden, saa følger heraf at det ene 0. ∞ maa kunne
være forskjelligt fra det andet, at med andre Ord 0.∞ er
en ubestemt Størrelse.

b) Den ubestemte Størrelse 0. ∞ kræver en Bestemmelse
af det uendeligt Lille og det uendeligt Store. Denne Bestemmelse beroer paa Forskjellen mellem faste og bevægelige
Grændser. Den faste Grændse kan ikke overskrides; den
bevægelige kan, saavel hvad Formindskelse som Formerelse
angaaer, overskrides i det Uendelige.

c) Da en Quotient af to foranderlige Størrelser forbliver
uforandret, naar Tæller og Nævner enten divideres eller
multipliceres med samme Størrelse, saa indsees heraf, at
Quotienter af uendelige (smaa eller store) Størrelser maae
kunne være ligesaa bestemte som af endelige.

Paa disse Forudsætninger beroer det Uendeliges Analyse,
hvis Udbytte for Philosophien er, at Modsigelserne mellem
det Endelige og det Uendelige, som ellers behandles i altfor
abstract Almindelighed, finde en punktuel dialektisk Løsning
paa et exact Grundlag.

% --- p. 2 ---
\setcounter{footnote}{0}%
\opage{2}\begin{center}\textbf{Propædeutik og Logik.}\end{center}


Philosophiens almene Grundlag er Logik. Denne Videnskab
kan behandles deels som formel Logik deels som Ontologie,
Metaphysik eller Ideernes Logik.

\begin{center}\textbf{Formel Logik}\\
\textbf{i}\\
\textbf{Grundrids.}\end{center}

Den formelle Logik eller Læren om de abstract nødvendige Tankeformer deler sig i Læren om Begreb, Dom og
Slutning.

\begin{center}\textbf{Begreb.}\end{center}


Tænkningen, som allerede er med i Sandsningen, hæver
sig frem, idet den opløser de umiddelbare Sandsebilleder, danner
ved Sammenligning og Abstraheren et Almeenbillede,
fastholder det i Ordet, udvikler dets nødvendigt sammenhørende
Væsensmærker og hæver det til et Begreb (\emph{Begrebsanalyse};
Definition).

Indbegrebet af de for et Begreb uundværlige Mærker
kaldes Begrebets \emph{Indhold}, Antallet af de Gjenstande, hvorpaa
det finder Anvendelse, kaldes dets \emph{Omfang} (Forhold
imellem Omfang og Indhold; over- og underordnede Begreber;
Slægts- og Artsbegreber).

\begin{center}\textbf{Dom.}\end{center}


Begrebet finder kun sit rene logiske Udtryk i Sproget;
Ord forbindes til Sætninger, og i Sætninger udtales
Domme. Den Sætning, der efter sin Form udtrykker et
Begreb, er en logisk Dom. For at udtrykke Begrebet, maa
Dommen formelig angive: \emph{Subjectet}, den Forestilling, der
skal bestemmes, \emph{Prædicatet}, den bestemmende Forestilling,
og \emph{Copula}, det bestemmende Forbindelsesord.

Af Subjectets ved Copula bestemte Forbindelse med Prædicatet
følger, at Dommen kan adskilles i en Fleerhed af lo%
% --- p. 3 ---
\setcounter{footnote}{0}%
\opage{3}gisk forskjellige Domme, og Dommene \emph{inddeles}: efter \emph{Qvalitet}
i bekræftende (a er b), benægtende (a er ikke b) og
uendelige (Samvittigheden er ingen Fiirkant); efter \emph{Qvantitet}
i singulære (dette a er b), particulære (nogle a ere b)
og generelle (alle a ere b); efter \emph{Relation} i kategoriske (a er,
ifølge sin Kategori, b), hypothetiske (dersom a er b, saa er
c d) og disjunctive (a er enten b eller c eller d); efter \emph{Modalitet}
i assertoriske (a er virkelig b), problematiske (a er
muligviis b) og apodiktiske (a er nødvendigviis b).

Fremgaaer Prædicatet med almeenlogisk Nødvendighed af
Subjectet, kaldes Dommen \emph{analytisk}, kan Prædicatet kun
ved en særegen Erfaring føies til Subjectet, kaldes Dommen
synthetisk (\textit{synthetiske Domme a posteriori} og \textit{a priori}).

\begin{center}\textbf{Slutning.}\end{center}


Dommen maa begrundes, dens Gyldighed indsees; den
Tænkningsact, hvorved en Doms Gyldighed udledes af een
eller flere andre Domme, er en logisk Slutning. Har Slutningen
kun een Forsætning (Præmis), kaldes den \emph{umiddelbar},
har den flere Præmisser, kaldes den \emph{middelbar} (Slutningens
Reqvisiter).

De \emph{middelbare} Slutninger \emph{inddeles} i: de \emph{kategoriske},
der betinges af, at det underordnede Begrebs Kjendemærke
(\textit{nota notæ}) subsumeres under det overordnede (syllogistiske
Figurer); de \emph{hypothetiske}, der beroe paa det Betingedes
Afhængighed af sin Betingelse (\textit{modus ponens} og \textit{modus
tollens}); og de \emph{disjunctive}, som grunde sig paa, at
Leddene af et Hele gjensidig forholde sig saaledes, at naar
eet udtrykkelig er sat, saa ere de øvrige udelukkede (\textit{modus
ponendo tollens}), og naar de øvrige ere udelukkede, saa maa
det tilbageblevne sættes (\textit{modus tollendo ponens}).

Ved at afhandle Læren om Beviis og Methode løser
den formelle Logik sig op i en almindelig Erkjendelses- og
Videnskabslære.

% --- p. 4 ---
\setcounter{footnote}{0}%
\opage{4}\begin{center}\textbf{Hovedsætninger}\\
\textbf{af Grundideernes Logik.}\end{center}

a) Den Gyldighed, der kan tillægges Philosophien,
beroer væsentlig paa den Gyldighed, der kan tillægges Ordet.
Det menneskelige Ord udtrykker kun en afmægtig Viden;
derfor falde Begreb og Gjenstand udenfor hinanden. Det
bliver i denne Henseende Philosophiens Opgave at paavise
Forholdet mellem det afledede og det oprindelige, mellem det
menneskelige og det guddommelige Ord.

b) Det guddommelige Ord er: α) i \emph{Videns} Form
det indre Billedlige, der bærer Forestillinger og Tanker, Begreber
og Domme, β) i \emph{Magtens} Form det udvortes Prægende,
den faktiske Udtalen af Begrebets og Gjenstandens
Eenhed; γ) i \emph{Vexelbestemmelsen} af Viden og Magt
den teleologiske Forskjelseenhed af Tænken og Virken, Dommen
og Handlen, Begriben og Skaben.

c) I \emph{Videns} Form bliver Forholdet mellem Dom og
Gjenstand opfattet theoretisk. Dommens Eenhed med Gjenstanden
forklares: α) ved en \emph{Parallelisme} (Plato, Cartesius)
saaledes, at Tingene opfattes som sandselige Afbilleder af
intellectuelle Urbilleder; β) ved \emph{Impression} og \emph{Abstraction}
(Hume) saaledes, at Begreberne blive intellectuelle Afbilleder
af de i Erfaringen givne Gjenstande; γ) ved en \emph{subjectiv
Determination} (Kant) saaledes, at de subjective
Erkjendelsesformer, der paavirkes af et ubekjendt objectivt
Værende (das Ding an sich), antages at bestemme Erfaringens
Phænomener.

d) I \emph{Magtens} Form bliver Forholdet mellem Dom og
Gjenstand at opfatte praktisk. Dommens Eenhed med Gjenstanden
fordrer da: α) en \emph{Præformation} (ideel Praxis),
hvorved Realbegrebet, det Begreb, der skal realiseres, forud
bestemmer Gjenstanden ved Henblik til de Betingelser, under
hvilke Realiseringen skal foregaae; β) en \emph{Realisation} (ideelreel
Praxis), hvorved det i Realbegrebet Almene Punkt for
Punkt omsættes i sit tilsvarende Individuelle, saa at Gjen%
% --- p. 5 ---
\setcounter{footnote}{0}%
\opage{5}standen, der frembringes, bliver at betragte som Begrebets
praktiske Definition; endelig γ) en Explication (Tilbagevenden
fra Praxis til Theorie), hvorved den Dømmende og Frembringende
viser sig at gjenkjende Frembringelsens Begreb i den
frembragte Gjenstand (Henblik til Skabelsen).

e) Den teleologiske Eenhed af Begrebet i \emph{Videns} med Begrebet
i \emph{Magtens} Form beroer paa interpreterende Functioner,
svarende til Vexelbestemmelser af \emph{at være og at betyde}.
Den Dialektik af Midler og Formaal, som herunder udvikler
sig, lader os indsee, at Virkeligheden udtrykker Begreber ved et
kategorisk Magtsprog.

f) Hvorledes det Subjective i \emph{Ideernes Logik} forholder sig
til det Objective, kan endnu yderligere oplyses ved den Udvikling
af Begreberne: Mulighed, Virkelighed og Nødvendighed,
der svarer til Dommenes Inddeling efter Modaliteten i den
formelle Logik.

\begin{center}\textbf{Propædeutik og Ideelære.}\end{center}


Ideen, Begrebets og Indholdets absolute Eenhed, --- det
i en Totalitet af Midler og Formaal udviklede Princip, Identiteten af det Subjective og det Objective --- kan siges at være
baade Eet med og forskjellig fra Virkeligheden. I Modsætning
til Phænomenet, det Ydre, er Ideen Væsenet, det Indre, og
dog er Ideen hævet over Adskillelsen af Væsen og Phænomen,
af Indre og Ydre, er i sin Modsætning til Phænomenet
Phænomenets og Væsenets væsentlige Eenhed, i sin Modsætning
til Virkelighedens Ydre den virkelige Eenhed af Ydre og Indre.

Da Ideen saaledes gjennemtrænger og bestemmer al
Virkelighed, saa følger heraf, at Ideelæren ikke kan faae Realitet
uden paa et omfattende Grundlag af logisk indledende Virkelighedsanalyser.
Vi have Exempler herpaa i Objectivitetens Logik, i
Læren om Functioner af Grunden, af Aarsagen og af Grundaarsagen.


\begin{center}\textbf{Hovedsætninger hørende til Ideelæren.}\end{center}


1) Grundaarsagen, Principet eller, som vi ogsaa kunde
kalde den, Idee-Aarsagen, udvikler sit Indhold i Totaliteter,
og \textbf{Totaliteternes Udvikling} gaaer igjennem Extremer og Slut%
% --- p. 6 ---
\setcounter{footnote}{0}%
\opage{6}ninger. Charakteristisk er i denne Henseende Forskjellen mellem
den \emph{mechaniske}, den \emph{chemisk-physiske} og \emph{organiske}
Totalitet, hvis logiske Formudvikling gjennem tilsvarende
Functioner angiver Forholdet mellem den \emph{kosmologiske},
den \emph{geologiske} og den \emph{biologiske} Idee og danner
Grundlaget for Naturphilosophien.

2) Modsætningen mellem Natur og Aand bestemmes i
Ideelæren saaledes, at Naturen, væsentlig seet, er Ideen i
Nødvendighedens Form, Aanden Ideen i Frihedens Form.
Natur og Aand ere i Ideen Eet; men denne Grundeenhed maa
dog opfattes saaledes, at begge i Kraft af Ideen tillige ere
hinanden oprindelig modsatte. Aanden kan ligesaa lidt afledes
af Naturen som Naturen af Aanden. Modsætningseenheden
er en Slutningseenhed, der kjendes paa den indre Uendelighed,
at Begyndelsen selv er en Slutning ligesaavel som Resultatet.

3) \emph{Aandsphilosophiens} Inddeling (efter Hegel) i de
tre Afsnit: a) Aanden i sit Begreb (Anthropologie, Psychologie);
b) \emph{den praktiske} Aand (Retsphilosophien, Historiens
Philosophie) og c) \emph{den absolute} Aand (Æsthetik, Religionsphilosophie)
har vel som Schema betragtet sine Fordele, men
maa dog med Hensyn paa Princip og Methode underkastes
en gjennemgaaende Kritik.
\seamOversigt

% --- p. 81 ---
\setcounter{footnote}{0}%
\opage{81}medens dog det abstracte Qvantum, der afgiver Qvantiteringens
umiddelbare Forudsætning, just skylder Qvantiteringen
sin Tilværen, og derved kommer til at forudsætte det, hvis
Forudsætning det selv skulde være: saa indsees heraf, at Qvantiteten
som abstract Begreb først kan faae sin sande Betydning
i og ved et Tilværende, hvis Qvantitet den er.

Om den rene Mathematik hedder det, at den finder
„Anvendelse“ paa virkelige Gjenstande (i Mechanik, Astronomie,
Optik o. desl.). Væsentlig forstaaet forudsætter en saadan
„Anvendelse“, at det Mathematiske oprindelig er virkeliggjort
i de Gjenstande, hvorpaa Mathematiken anvendes. En Anvendelse
af Mathematik paa virkelige Gjenstande vilde være,
om ikke en reen Umulighed, saa dog en subjectiv Vilkaarlighed,
dersom det Mathematiske ikke selv med indre objectiv
Nødvendighed hørte med i Virkelighedens Væsen; men som
Moment i Virkelighedens Væsen hører det Mathematiske
tillige med i Philosophien.

\begin{center}\rule{4em}{0.4pt}\\[0.5em]
B.\\[0.3em]
\textbf{Propædeutik og Ideernes Logik.}\\[0.5em]
\rule{4em}{0.4pt}\end{center}


For at erkjende den Betydning, Ordet som Ord maa
udøve paa Grundforestillingerne om det Evige, paa Ideernes
første naturlige Dannelse i Bevidstheden mærke vi os særlig
Kategoriernes Vorden i Mythologie og Theosophie.

Kjendemærket paa det Mythiske er, i Principet, ingenlunde
Uklarhed over det udvortes Endelige; det er ikke Sandselighedens
Taagebilleder, men Ideens Lynslag, ikke Forstandens
Nat, men Fornuftens Daggry, der laaner Bevidsthedsindholdet
de vidunderlige mythiske Farver. Hvor eiendommelig Farveblandingen
er, falder i Øinene, naar man ret vil charakterisere
Mythologien og det mythiske Standpunkt. „For Mythologien“,
% --- p. 82 ---
\setcounter{footnote}{0}%
\opage{82}siger Heiberg, „gives der intet andet, Alt omfattende Synspunkt,
end det poetiske. Mythologien er Poesie og reent
oprindelig Folkepoesie“. Senere hen bliver saa denne Folkepoesie
tillige opfattet som en Art speculativ Philosophie.
„Enhver Mythologie“ hedder det, „begynder med Kosmogonie;
thi den philosopherende Menneskeheds første Bestræbelse er at
forklare Tilværelsens Under. Det første Spørgsmaal er
hvorledes er Verden blevet til? Hvad var der førend Alt
dette blev til? Eller hvad var der, da der endnu Intet var?
Forstanden kan saalidet besvare disse Spørgsmaal, at den
endogsaa maa erklære dem for taabelige og sig selv modsigende.
Men Phantasien svæver dristigt hen over Tingenes Begyndelse,
uden at frygte for at tabe sig i Intet. Den befolker Tomheden,
den lader Lyset opstaae af Mørket, Tilværelsen af
Ikketilværelsen“\footnote[1]{J. L. Heibergs nordiske Mythologie. Paa Dansk udgivet af Chr. Winther. Kbhvn. 1864. S. 4 og S. 24.} o. s. v.

At Mythologien paa een Gang er baade Poesie og
Philosophie, vil da med andre Ord sige, at den i Grunden
hverken er Poesie eller Philosophie; ikke Poesie, fordi den er
Philosophie; ikke Philosophie, fordi den er Poesie. Det
Nærmere sees af Functionerne. Det Poetiske bliver her ikke
opfattet som Poesie. Poesien forudsætter en paa det absolute
Indhold lydende Adskillelse af Illusion og Virkelighed, en
Adskillelse, der udtrykkelig sætter Illusionen over Virkeligheden,
fordi den sætter Ideen over Endeligheden; men ved en saa
reflecteret Adskillelse kan Mythen som Mythe umulig bestaae:
de mythiske Skikkelser holde sig ved Ideens Magie paa en
høiere Virkeligheds Trin; her er paa dette Trin end ikke
Anelse om Illusion. Ligesaalidt bliver det Philosophiske
i Mythedannelsen opfattet som Philosophie. Hvad skulde det
vel være for en Art Philosophie, som begyndte med en
udtrykkelig Splid mellem Forstand og Phantasie, og som paa
Grund af denne Splid opgav Forstanden og lod Phantasien
% --- p. 83 ---
\setcounter{footnote}{0}%
\opage{83}raade? Istedetfor at begynde med en Dualisme af Phantasie
og Forstand, begynder Mythedannelsen tvertimod med en af
Forstand og Fornuft imprægneret Phantasie; man kan om
den mythiske Ideebevægelse med fuld Ret sige, at Tanken
phantaserer, fordi \emph{Phantasien tænker.}

Naar da den af Forstand og Fornuft imprægnerede
Phantasie „svæver dristigt hen over Tingenes Begyndelse,
uden at frygte for at tabe sig i Intet“, hvad styrer da dens
Flugt, hvad er det, der bevarer den fra at styrte i Tomhedens
Afgrund? Væsentlig er det \emph{Ideen}; umiddelbart seet, er det
\emph{Ordet.} Begyndelsens Afgrund er lige dyb, enten vi begynde
med Chaos eller med \emph{Ginungagap}, den er i begge uendelig
dyb: men i det uendelige Dyb formæles Ideen med Ordet.

Begynde vi med Chaos, da er det naturligviis ikke den
tomme Benævnelse, hvorpaa her lægges Vægt, thi Benævnelserne
vexle. Men de eenstydige Benævnelser lyde alle paa det samme
Grundord. At Chaos nu som Grundord er en Ideens Aabenbarelse,
sat og bevæget af Ideens Magt, kan let paavises.
Chaos er i mythisk Forstand hverken et abstract logisk Begreb
eller et Udtryk for en Tilstand af physisk Forvisning, men
en i Ordet fastholdt Totalitetseenhed, et subjectivt-objectivt
Første, der i billedlig og dog ubilledlig Form anskueliggjør, hvad
ikke kan anskues, nemlig Tingenes Oprindelse, Ideen af al
virkelig Begyndelse eller den virkelige Begyndelses Idee.

At Chaos som Begreb ikke er i Begrebets, men i et
ubilledligt Billedes Form, følger udtrykkelig af Anskuelsens
Eenhed med Ordet. Hvad er der i Chaos? Intet, og dog Alt!
Denne Tankens Selvmodsigelse med de deraf udspringende
Conseqvenser opfattes ikke i Form af en Tanke, men sees
i et Billede, et mægtigt Tankebillede, dannet af Ordet.
Idet Tanken ved Ordets Hjælp phantaserer over Intet og
Alt, tænker Phantasien væsentlig paa Alt og anskuer dog
i Chaos Intet, det Tomme; den tænker væsentlig paa Intet,
og anskuer dog i Chaos Alt, nemlig Spirerne til Alt (\textit{non
bene junctarum discordia semina rerum}). Indholdet,
% --- p. 84 ---
\setcounter{footnote}{0}%
\opage{84}det i Chaos indholdsløse Indhold, er væsentlig subjectivt
formedelst Ordet: et Product af phantaserende Tænkning; af
tænkende Phantasie; men det anskues tillige objectivt, slaaer
den anskuende Bevidsthed med Objectivitetens Magt i Ordet,
og antager Skikkelse af physisk umiddelbar Realitet.

Men at Chaos i objectiv Mening desuagtet ikke har
physisk umiddelbar Realitet, er netop en Følge af det i
Billedet Ubilledlige, af Anskuelsens mythiske Eenhed med Ordet.
I objectiv physisk Betydning skulde Chaos jo være den virkelige
Opløsning, hvori Altilværelsen maatte befinde sig, naar alle
Former vare forstyrrede, og Atomernes Vrimmel i regelløs
Forvirring. For en forstandig kritisk Bevidsthed om et
saadant Chaos har Ordet tabt sin magiske Kraft. Ordet
Chaos er i dette Tilfælde kun et Navn; Chaos derimod som
udvortes virkelig Tilstand er Tingen, og Tingen, ikke Ordet,
er det, hvorpaa det kommer an. Adskilt fra den paa Ordet
hvilende Phantasie, vil Forstanden da paa egen Haand
gjennemforske det virkelige Chaos, Tingenes udvortes virkelige
Begyndelse. Men dette Forstandens Forehavende vil totalt
mislykkes. Som virkelig Tilstand er Chaos aldeles uvirkelig;
der er virkelig ingen Tilværelse, der kan begynde med et
virkeligt Chaos: hvad Forstanden skal faae ud af et saadant
Chaos, maa den i Virkeligheden selv indsmugle deri: det er
en ørkesløs Leg, Forstanden leger med sig selv og det Virkelige.

I det mythiske Chaos er Objectiviteten subjectiv; det
objective Indhold er ved Ordet omsat i subjectiv Form.
Hvor forskjelligt det mythiske Chaos er fra det udvortes
physiske, er da ogsaa kjendeligt paa de Skikkelser, der fremgaae
af Chaos. „Først er Chaos, dernæst den udstrakte Jord,
det mørke Tartaros --- og Eros, den skjønneste blandt de
udødelige Guder“. Den, der seer de gamle Guder udtræde
af Chaos, kan ikke være i Tvivl om, hvad Chaos er for et
Væsen. Idet her begyndes med Chaos, tages her paa en
Maade Tilløb til en objectiv Fremstilling af Verdens
Tilblivelse, af Tingenes, de endelige Realiteters Dannelse;
% --- p. 85 ---
\setcounter{footnote}{0}%
\opage{85}men Synet vender sig. Kosmogonien slaaer om og bliver
Theogonie. Chaos bliver en Guddom, den Guddom, af hvis
Skjød de særlige Guder fødes. Denne Tilbagegang af det
Kosmogoniske i det Theogoniske er for Mythedannelsen i høi
Grad charakteristisk. For at Verden kan blive til, maae
Guderne først blive til; for at de endelige Realiteter i
Virkeligheden kunne blive til, maae Realideerne først blive til:
Realideernes endelige Vorden reflecterer i sig Grundideernes
evige Væren.

Chaos er altsaa i sin Idealitet evig, i sin umiddelbare
Realitet timelig, Guddom i Væren, Naturgrund i Vorden,
Totalitet i oprindelig Eenhed, Totaliteternes Grund i oprindelig
Fleerhed; i Chaos er det Relative absolut, det Absolute relativt;
i Chaos-Ordet er Bevægelsen Hvile, og Hvilen igjen Bevægelse;
Fylden Tomhed, og Tomheden Fylde; det Indvortes et Udvortes,
og dette Udvortes igjen et Indvortes: i Chaos er Ideen Ord,
og Ordet Idee.

Det er da ikke en Tilfældighed, men en væsentlig
Nødvendighed, at de ældste Guder udgaae af Chaos. I logisk
Form er Nødvendigheden kjendelig paa, at Fleerheden fremgaaer
af Eenheden, efterdi Eenheden selv oprindelig er
Fleerhed, at det Særskilte i Totaliteten fremgaaer af det
Almene, efterdi det Almene oprindelig er en Totalitet af
Særskilte. Men i den logiske Form mangler Indholdet
umiddelbar Realitet; Realiteten, den mythiske Realitet,
kommer med Ordet, thi det ved Ideen bestemte Ord bestemmer
Ideen.

Saadanne Skikkelser som Gæa, Tartaros, Erebos, Nyx
og Eros ere i deres Oprindelighed mere end tomme Personificationer.
„Skikkelser af det Skikkelsesløse“, „elementære
Former af det første Ene, i sig selv Formløse“, slige fastende
Abstractioner falde lige saa langt under Ideen som under
Virkeligheden. Nei, ved Ideens Magie i det mythiske Ord
ere de ældste Guder, trods det i Formen Uformelige, typiske
Skikkelser, fulde af Realitet, hver Guddom en Totalitet i sig,
% --- p. 86 ---
\setcounter{footnote}{0}%
\opage{86}et eget Væsen, et bestemt Physiognomie. Det er Ideen, der
i særskilte Totalitetseenheder gjennemtrænger, bevæger og bestemmer
Ordet; men det er i umiddelbar Forstand Ordet,
det særskilte Grundord, der i sin Eenhed optager, bestemmer
og samler Mangfoldigheden af de ideelle og reelle, ubilledlige
og billedlige Momenter, der udgjøre Ideens absolute Indhold.

Gudinden Gæa, der faaer Navn efter Jorden, optager
ved dette Navn al Jordens Realitet i sig, bliver ved dette
Navn en Guddom med umaadelige Evner, vældige Kræfter
og store Muligheder; men at Gudinden har faaet Navn efter
den virkelige Jord, siger med det Samme, at det er den
virkelige Jord, der har faaet Navn efter Gudinden. Jordgudinden
og den virkelige Jord have ikke blot eet Navn, men
de ere eet Væsen, betegnede ved eet Ord: de ere i Ordet Eet,
efterdi de ere i Ideen Eet. At den virkelige Jord er Gudinde,
viser os ret paa en slaaende Maade, hvad det vil sige, at
Ideens objective, reelle Indhold er --- formedelst Ordet
--- omsat i subjectiv, ideel Form; at Gudinden med alle Træk
af elementær Subjectivitet er Jorden selv, den virkelige Jord,
anskueliggjør paa den anden Side, hvad Ordet formaaer,
naar det gjælder om at oversætte et subjectivt Indhold, ---
thi Gudindens Væsen er i sin Idealitet et subjectivt Væsen, ---
i umiddelbar, objectiv Form.

Den colossale Dristighed, hvormed den tænkende, af
Ideen imprægnerede Phantasie ved Syn af virkende Guder
trodser Sandsynligheden, er et talende Vidnesbyrd om Ordets
Magt i de objectiverede Begivenheder. Gudealderen, den
fjerne Fortid, er Ideernes, de ideale Aarsagers gyldne Tid;
Nærværelsen med sine Kjendsgjerninger er Virkelighedens, Betingelsernes
og de materielle Aarsagers Tid: kunne de to
Tider vel blive til een Tid? I Ordets, d. v. s. i Fortællingens,
Erindringens, Speil reflecterer det Forbigangne sig
jo vistnok i det Nærværende som de store Grundideer i Udviklingens virkelige Resultater; men Virkeligheden med sine
Tingsaarsager reagerer med Nærværelsens Magt mod hine
% --- p. 87 ---
\setcounter{footnote}{0}%
\opage{87}ideale Forudsætninger; Conseqvensen er en Splid mellem
Tingen og Ordet, en Splid mellem det Objective og det
Subjective, en Splid mellem Indhold og Form, Phantasie
og Tanke, en Splid, der vækker Kritiken, ophæver Standpunktet
og stiller de ideeformende Functioner nye og høiere
Opgaver.

Saavidt om det Mythiske. Til Belysning af den
theosophiske Form kaste vi et Blik paa Philo.

„Naar man“, siger en nyere Forfatter, „beskjæftiger sig
med den jødiske Alexandrinisme, henvises man uvilkaarligt til
Hebraismen og Platonismen, som dem, under hvis Paavirkning
den uimodsigeligt har uddannet sig. Heri ere alle
Fortolkere enige; Forskjellen fremtræder først, naar man skal
gjøre sig det klart, hvilken Aandscharacteer, der udtalte sig
deels i selve denne Receptivitet, deels i den friere Productivitet,
som denne kaldte tillive. I Gfrører og Dähne finde vi nu
de seneste og meest udprægede Repræsentanter for Opfattelsen
af den som reflecterende; thi skjøndt de polemisere skarpt mod
hinanden og ende i saa godt som modsatte Resultater, fremstille
de dog det Standpunkt, hvorpaa de tænke sig Philo
saavelsom hans Forgjængere i den alexandrinske Theosophi
at staae, hovedsageligt eens. --- Som de aandelige Kræfter,
hvis Sammenstød bringer Bevægelsen ind i Philos Aand,
betragte de fra den ene Side den bestemte Tilegnelse af den
platoniske Modsætning mellem det Guddommelige og Materien
og af den deraf med Nødvendighed følgende Tanke, at Gud
ikke kunde staae, end ikke i den fjerneste Forbindelse med
Endeligheden, saa at endog alle Prædikater, som hentede
fra denne, bleve at frakjende ham; fra den anden Side
den ligesaa uopgivelige Tro paa Gud som Person, paa
hans Ledelse af Jødefolket og Moselovens Guddommelighed.
Den Aandsvirksomhed, som derved tænkes fremkaldt, var da
en Bestræbelse efter at forbinde disse to modsatte Elementer
til eet System, og til den Ende skal der være uddannet den
dobbelte Theori: om den allegoriske Skriftfortolkning og de
% --- p. 88 ---
\setcounter{footnote}{0}%
\opage{88}guddommelige Mellemvæsener. Som den, der skal være
baaren af denne bevidste Stræben efter at lade begge disse
to modsatte Principer holde hinanden Stangen, saa at begges
Hovedsætninger bleve urokkede, characteriseres ogsaa Philos
Lære som en rationel Supranaturalisme.

Naar vi nu i Et og Alt fjerne os fra denne Gudsbetragtning,
saa er det ikke nærmest fordi den skulde støde an
imod et enkelt Punkt, men fordi den modstrider hele den
philoniske Erkjendelses Characteer. Er den nemlig fremstaaet
ved slig combinerende Reflexion, er dens enkelte Forestillinger
og Udsagn Udtryk for en saadan, saa maa den ogsaa bære
dennes Kjendetegn. Men det er det, den ikke gjør. Ad en
dobbelt Vei kunde nu denne Reflexion gaae frem, hvorefter
vel Philos Værker vilde faaet et høist forskjelligt Præg, men
dog altid baaret Reflexionens Characteer. Man skulde snarest
have troet, at, hvis Philo virkelig havde havt den Vanskelighed
for Øie, som kom frem ved de før berørte modsatte Fordringer,
men dog betragtet den som en, der ved nøiere Undersøgelse
forsvandt, thi var den forekommen ham uopløselig, maatte
han jo have bestridt en saadan Sammensmeltning og ikke selv
forsøgt den, --- hans Bestræbelse da nærmest maatte være
gaaet ud paa, ved raisonnerende Sammenstilling at eftervise,
hvorledes denne Modsigelse kun var tilsyneladende og hævedes
--- vi vide jo hvorved --- ved Allegorien og ved Forestillingen
om de guddommelige Mellemvæsener; og hans Philosophi
vilde paa denne Maade have taget en \emph{formel-logisk} Retning.
\emph{Men til alt dette findes hos Philo aldeles intet
Spor.} Hvad nemlig \emph{Allegorien} angaaer, som jo skal være det
ene Mellemled, da vil man vel . . . . kunne slutte, hvorfor vi paa
ingen Maade kunne dele denne Mening, men ikke desto mindre
ville vi dog rette en Bemærkning imod den. Skal Allegorien
være uddannet i dette Øiemed, nødes man jo til at
tænke sig, at Philo, med Øiet fæstet paa den omtalte
Vanskelighed, med Ønsket om at løse den, hvis det var
muligt, uddannede denne Theori, ikke som et Middel til at
% --- p. 89 ---
\setcounter{footnote}{0}%
\opage{89}blænde Taaber, men som det, der tilfredsstillede ham selv,
idet den ikke syntes ham fremtvungen ved et vilkaarligt Magtsprog,
men i sig selv at være sand, thi først da kunde den
jo atter bevise Sandheden af det Øiemed, hvori den brugtes“.
Men saaledes forholder det sig ikke.

At det ikke er „rationel Supranaturalisme“, men
„Theosophie“, hvorom her handles, kan sees af Forholdet
mellem Ordet og Gud.
% p. 89 from here: the setting of GL II p. 89; the draft takes texts/nielsen/grundideernes-logik-2 (same type), then collated

1) \emph{Ordet og Gud i tvetydig Eenhed}. Den philoniske
Hovedkategorie er, som bekjendt, Kategorien Logos; men
om Philo iøvrigt opfatter Logos som identisk med Gud eller
som et fra Gud forskjelligt Væsen, om dette Væsen skal antages
at være personligt eller upersonligt, derom have Fortolkerne
hidtil ikke kunnet blive enige\footnote[1]{Jvfr.\ G.\ Steenberg. Om Synspunktet for Opfattelsen af Philos
Gudserkjendelse. Kjøbenhavn 1849. S.\ 84---125.}. Den theosophiske
Reflexion vælger i al Naivetet forskjellige Udgangspunkter,
uden at ændse de Modsigelser, som fremkomme, naar
ueensartede Synsmaader skulle gaae op i hinanden. Paa
den ene Side fremstilles da Logos som et med Gud identisk
Væsen, som Guds Tænkning (\gr{λογισμὸς τοῦ θεοῦ}), som Guds
Ord (\gr{ῥῆμα}), Guds Eed (\gr{ὅρκος}), Guds Lov (\gr{νόμος}), Guds
Pagt (\gr{διαθήκη}). Men ligesaa umiddelbart opstilles Logos
paa den anden Side som et Væsen for sig, som det Baand
(\gr{δεσμὸς}), der sammenbinder alle Ting (\gr{περιέχων τὰ πάντα}),
som, „for at hindre Alt fra at skeie ud, tænkes udspændt fra
den ene Ende af Verden til den anden“; Logos er den guddommelige
Viisdom, der gjennemstrømmer Verden; Verden % PRINTER (GL II): quotation opened at „den guddommelige is not closed
er denne Viisdoms Klædebaand; Viisdommen selv er \gr{ὁ τῆς
φύσεως ὀρθὸς λόγος}. Uagtet Logos i begge disse Opfattelser
er upersonlig, saa indstille sig dog igjen andre Synspunkter,
fra hvilke Logos udtrykkelig sees som personligt Væsen,
som en Gud selv underordnet Gud (\gr{δεύτερος θεὸς, πρωτόγονος
υἱὸς}), o.\ s.\ v.

\seamB

% --- p. 263 ---
\setcounter{footnote}{0}%
\opage{263}enkelte Sætninger derimod ikke særlig fremtræde med apodiktisk
Selvstændighed, da hver relativ Bestemmelse er et Moment,
et Glimt af det Absolute. Men just denne de apodiktiske
Momenters reflecterede og i Reflexionen gjennemsigtige Form
gjør Ideen, det absolut Apodiktiske, synlig.

% p. 263: the setting of GL II p. 263 + its p. 264 line 1; the draft takes texts/nielsen/grundideernes-logik-2, then collated.
% The note at the foot is the end of a note begun on p. 262 (GL II p. 262, at „Uendelighedsproblem.“).
Modalitetens almindelige Hovedformer: Mulighed, Virkelighed
og Nødvendighed, anbefale sig vexelviis til den sunde
Sands saaledes, at snart Endeligheden, snart Uendeligheden,
snart Virkelighedssiden, snart Ideesiden vender imod Betragteren.
Paa Endelighedens Trin gjælder det, at der nødvendigviis maa være Noget, der blot er \emph{muligt}, Noget, der er
\emph{virkeligt}, og Noget, der er \emph{nødvendigt}. Seet under
Uendelighedens Synspunkt træder Forskjellen mellem Muligt,
Virkeligt og Nødvendigt saaledes i Baggrunden, at den ene
af Formerne uvilkaarligt faaer Magt med de to andre, og nu
\contnote{have kunnet faaet det i deres Tanker ubegribelige Begreb (Uendehedsbegrebet) paa en tilfredsstillende Maade adskilt fra den paa % PRINTED AS IS: „(Uende-/hedsbegrebet)“, a syllable lost (crop k2); GL II 1866 has „(Uendelig-/hedsbegrebet)“ and „faae“ where 1868 has „faaet“: the first line of this note was reset for 1868
dette Begreb hvilende Methode, viser Mathematikens Historie lige
fra Leibnitz og Newton (Jvfr.\ Mathem.\ og Dialekt.\ S.~90--91,
S.~95 o.\ flg.) indtil den nærværende Tid. „For det Første kunne
vi, naar Principet er ubekjendt, jo ikke vide, om ogsaa alle deraf
følgende Conseqvenser lade sig indrømme: vi gaae som i Blinde,
og det turde hænde sig, at vi efter møisommelig Anstrengelse tilsidst
stødte paa Noget, som netop udviste Principets Natur, d.\ v.\ s.
paa Noget, vi ikke kunde begribe. For det Andet: er det ikke en
skammelig Nedværdigelse af Fornuften, at lade den følge en Sammenhæng
af Led, en Række, hvis Begyndelse og Ende den ikke kan
gjennemskue?“ Disse Savriens Betænkeligheder forklare tilstrækkelig
den Uro og Spænding, hvori det Uendeliges Analyse bestandig har
holdt den mathematiske Reflexion (Mathem.\ Dial.\ S.~55--56),
hvis Problem --- saaledes som en af Berlinerakademiet i Aaret
1784 udsat Priisopgave charakteristisk nok udtrykker det, „wie aus
einer widersprechenden Voraussetzung so viele wahren Sätze haben
können hergeleitet werden“ --- ikke vil ophøre at vende tilbage.
Forelæsninger over „Phil. Propæd.“ fra Universitetsaaret 1860--61.
S.~47--48.}
% --- p. 264 ---
\setcounter{footnote}{0}%
\opage{264}hedder det: snart, at Alt er muligt, snart, at Alt er virkeligt,
snart, at Alt er nødvendigt. Falder Eftertrykket paa Vorden
(Heraklit), saa er Alt muligt: falder Eftertrykket paa Væren
og det Værende, saa er Alt virkeligt (Megarikerne); falder
Eftertrykket paa Virkelighedens i sig reflecterede Væsen, den
uendelige Substans, saa er Alt nødvendigt (Spinoza). I logisk
Forstand er Nødvendigheden overgribende, men den overgribende
Nødvendighed er subjectiv væsentlig.

\begin{center}\rule{4em}{0.4pt}\\[0.5em]
C.\\[0.3em]
\textbf{Propædeutik og Ideelære.}\\[0.5em]
\rule{4em}{0.4pt}\end{center}


Naar Ideen skal erkjendes som Virkelighedens Idee og
optage Kjendsgjerningsmomenterne i sit væsentlige Indhold,
vil Udviklingen fornemmelig være betinget af en indtrængende
Analyse ikke blot af det Objectives Eenhed med det Subjective,
men af Objectivitetsbestemmelsernes egen uendelige Rigdom.
Saalænge det Objective blot behandles i al Almindelighed,
maa dets eiendommelige Charakteer nødvendig forflygtiges;
thi det Almindelige er ikke, hvad man i Modsætning til
det psychologisk Subjective altfor eensidig har fremhævet,
det Objective, men tvertimod det i sit Væsen oprindelig
Subjective: at behandle det Objective i al Almindelighed
er at behandle det paa en subjectiv Maade, er, med
andre Ord, at behandle det, ikke som et Objectivt, men
som et Subjectivt. Sættes f. Ex. Idee og Virkelighed imod
hinanden, saa træder Ideen, uagtet den selv er Eenheden af
det Subjective og det Objective, dog øieblikkelig over paa den
subjective Side, medens Virkeligheden bliver staaende som særligt
Udtryk for det Objective, Grunden hertil er let at indsee.
Kjendemærket for det Objective er overalt Andetheden, Enkeltheden
i dens Umiddelbarhed, det udvortes Mangfoldige, Forskjelligartede o. s. v. Nu kan der jo vistnok i al Almindelighed siges

\seamC

\end{document}
